\section{Bilan Professionnel et Technique}

Ce stage au sein de TamTam m'a certainement permis de connaître et sentir la différence entre la vie d'un étudiant et la vie professionnelle d'un ingénieur dans le monde de l'entreprise. Être intégrer dans une équipe d'ingénieurs avec plus d'expérience, et le fait de travailler sur un projet réel, dans laquelle, toute fonctionnalités ajoutées il les fallait bien tester avant de passer en prod, m'a fait sentir beaucoup de responsabilités.

Il a fallu prendre la parole, avoir un esprit de synthèse, défendre mes arguments, être prêt à toute éventualité, être rapide dans la réalisation des tâches qui m'ont été accordées sans oublier la qualité, j'ai senti que je suis un salarié non plus un stagiaire.

Une bonne communication et s'entendre bien avec tous les membres de l'équipe étaient aussi des grands facteurs dans la réussite de mon travail.

De plus, ce stage m'a permis de découvrir et travailler avec plusieurs technologies à la fois, que je n'avais jamais eu l'occasion de les apprendre, ainsi que pouvoir travailler sous plusieurs normes et principes, cette excellente organisation avait permis de réaliser un produit final très bien fait.

En effet, travaillé sous la responsabilité d'un ingénieur logiciel bien compétent m'a offert l'opportunité de bien maîtriser les bonnes pratiques de travail au niveau d'organisation du code (clean code), de conception, d'optimisation, de tests et aussi l'utilisation des différents outils qui assurent le bon déroulement des projet (Git, Notion, Symfony…).

\section{Bilan Personnel}

Personnellement, ce stage m'a permis de vivre des grands défis. Il fallait être à jour, avoir un esprit d'initiative, renforcer mes relations d'amitié.

Vraiment, je suis totalement fier du travail effectué durant cette période. Surtout que le projet est mis en production et le processus de synchronisation est en cours.

\section{Apprentissages et Compétences Acquises}

\subsection{Compétences Techniques}

\begin{itemize}
    \item Maîtrise approfondie de Symfony et du framework PHP
    \item Développement en React.js avec Redux pour la gestion d'état
    \item Connaissance des architectures modernes (Hexagonale, CQRS)
    \item Implémentation de systèmes de synchronisation de données
    \item Utilisation de Git et GitHub pour le contrôle de version
    \item Gestion des bases de données MySQL avec Doctrine ORM
    \item Déploiement avec Nginx et GitHub Actions
end{itemize}

\subsection{Compétences Méthodologiques}

\begin{itemize}
    \item Application pratique de la méthodologie SCRUM
    \item Participation aux sprints et daily standups
    \item Utilisation de Notion pour le suivi des tâches
    \item Bonnes pratiques de développement (Clean Code, Tests)
    \item Documentation et communication technique
end{itemize}

\subsection{Compétences Professionnelles}

\begin{itemize}
    \item Communication efficace en équipe
    \item Travail collaboratif dans un environnement professionnel
    \item Gestion du temps et des priorités
    \item Adaptabilité face aux changements
    \item Prise de responsabilité et d'initiative
end{itemize}

\subsection{Compétences Personnelles}

\begin{itemize}
    \item Autonomie et capacité à apprendre rapidement
    \item Persévérance dans la résolution de problèmes
    \item Confiance en soi et estime de soi renforcées
    \item Développement du leadership et de la prise de décision
end{itemize}

\section{Conclusion}

Ce stage de fin d'études a été une expérience enrichissante et déterminante pour ma carrière professionnelle. Il m'a permis de transformer les connaissances théoriques en applications pratiques et de développer des compétences essentielles pour ma future carrière en tant qu'ingénieur informatique.

La réalisation du système de synchronisation incrémentale pour PowerTeam m'a confronté à des défis techniques complexes et m'a permis de démontrer ma capacité à les surmonter. Le travail en équipe, l'adoption de la méthodologie Agile et la mise en pratique des architectures modernes ont contribué à ma croissance professionnelle.

Je suis reconnaissant pour cette opportunité et pour le soutien de l'équipe de Tamtam International qui m'a permis de me développer à la fois sur le plan technique et personnel.
